\documentclass[level=2,teacher]{kidscode}
% Let lines with long code words stretch a little instead of running into the margin.
\setlength{\emergencystretch}{3em}
\begin{document}
\kctitle{Teacher Guide}{Everything you need to run Level 2 at home: setup, the weekly routine, assessment, pacing and troubleshooting}

\begin{teachernote}[Read this first (5 minutes)]
Welcome back! In Level 1 your group moved from Scratch blocks to Python turtle. Level 2
keeps that same friendly weekly rhythm, but now the children write \emph{real} programs:
programs with their own functions, lists of things, dictionaries, files that remember
high scores, automatic tests, and arcade games that react to keys.

You do \textbf{not} need to be a programmer. Each lesson plan tells you what to say,
what to click and what to look for, and every code answer in the Answer Key is a file
that has been tested by a computer. Your job is to keep the session calm and curious,
to keep time, and to ask good questions. When a child is stuck, \enquote{Let's read the
error message together} is almost always the right first move.
\end{teachernote}

\setcounter{tocdepth}{1}
\tableofcontents
\clearpage

%% ====================================================================
\section{1. Welcome to Level 2: what changed since Level 1}

\begin{tabularx}{\linewidth}{@{}>{\bfseries\raggedright\arraybackslash}p{3.4cm}X@{}}
\toprule
In Level 1 & In Level 2 \\ \midrule
Scratch, then Python & Python in Thonny every week. No Scratch. \\
Short scripts & Programs organized into functions with \texttt{return}, with the main program under
\texttt{if \_\_name\_\_ == "\_\_main\_\_":} (the lesson plans explain this line when it first appears). \\
Variables & Values of different \textbf{types} (whole numbers, decimals, text, True/False), lists and dictionaries. \\
Programs forget everything & Programs \textbf{save and load files} (week 7). \\
\enquote{Does it look right?} & \textbf{Automatic tests} with pytest from week 8. Every exercise folder has a
\texttt{check\_my\_work.py} the children can run themselves. \\
Drawing pictures & Event-driven \textbf{games} (keys and timers) and generative art. \\
Driver and navigator & Driver, navigator \textbf{and code reviewer} (section~6). \\
Homework 10--15 min & Homework 15--20 min: Core, plus optional Challenge and Creative tiers. \\
No placement quiz & A 15-minute \textbf{placement quiz} at the start of week 1 (section~8). \\
\bottomrule
\end{tabularx}

\begin{tip}[What stays the same]
The five-part session (Warm-up, Concept, Guided demo, Hands-on build, Wrap-up), the weekly
bug hunt, unplugged activities, pair programming, creative free-build, reflection, and
the digital-citizenship thread. Children who loved the Level 1 routine will feel at home.
\end{tip}

\subsection{The documents for Level 2}
\begin{tabularx}{\linewidth}{@{}>{\bfseries}p{4.2cm}X>{\raggedright\arraybackslash}p{2.6cm}@{}}
\toprule
Document & What it is for & Who \\ \midrule
Teacher Guide & This book: the big picture, setup, assessment, pacing, troubleshooting. & You \\
Lesson Plans & One plan per session (weeks 1--12, buffers B1 and B2). Read one the day before. & You \\
Student Workbook & One chapter per week, used during the session. & Each child \\
Homework & One 15--20 minute set per week. & Each child \\
Answer Key & Answers to the workbook, homework and assessments. & You (keep it away from the children) \\
Assessments & Placement quiz, week-6 checkpoint, week-12 exit assessment. & Each child \\
Project Brief / Project Teacher Guide & The final project (weeks 9--12). & Child / You \\
Extras & Cheat sheet, glossary, progress tracker, parent notes, certificate. & Everyone \\
\bottomrule
\end{tabularx}

%% ====================================================================
\section{2. Level 2 goals}
By the end of Level 2 each child can do the things below. The codes (L2.G1 \ldots) appear
in every lesson plan next to the learning objectives, so you can see which goal a lesson
builds. The CSTA codes refer to the CSTA K--12 Computer Science Standards (2017), a
widely used list of what children should learn in computing; you do not need to look them up.

\begin{longtable}{@{}>{\bfseries}p{1.4cm}p{9.6cm}>{\raggedright\arraybackslash}p{3.6cm}@{}}
\toprule
Code & By the end of Level 2 the student can\ldots & CSTA \\ \midrule\endhead
L2.G1 & use clearly named variables of different types (int, float, str, bool) and f-strings & 2-AP-11 \\
L2.G2 & combine control structures: while loops, nested loops, compound conditions & 2-AP-12 \\
L2.G3 & write functions with parameters and return values to organize code & 2-AP-14, 2-AP-13 \\
L2.G4 & store and process collections with lists and dictionaries & 2-AP-11, 3A-AP-14 \\
L2.G5 & read and write text files to save and load data & 2-DA-08 \\
L2.G6 & debug systematically (read tracebacks, print-debugging, rubber duck) and write simple tests & 2-AP-17, 2-AP-19 \\
L2.G7 & build event-driven games and generative art with turtle & 2-AP-12, 2-AP-16 \\
L2.G8 & plan with pseudocode/flowcharts, work in a team and act on feedback & 2-AP-10, 2-AP-15, 2-AP-18 \\
L2.G9 & make safe, fair, private choices about data and media & 2-NI-05, 2-NI-06, 2-IC-21, 2-IC-23 \\
\bottomrule
\end{longtable}

%% ====================================================================
\section{3. The 12-week roadmap}
Weeks 1--5 bring new ideas, week 6 is review and the mid-level checkpoint, weeks 7--9 bring
more new ideas, weeks 10--11 are the final project studio, and week 12 is the showcase and
exit assessment. The project is introduced at the end of week 9.

{\small
\begin{longtable}{@{}>{\bfseries}p{0.8cm}>{\raggedright\arraybackslash}p{3.2cm}>{\raggedright\arraybackslash}p{5.6cm}>{\raggedright\arraybackslash}p{5.2cm}@{}}
\toprule
Wk & Title & Big idea & New concepts \\ \midrule\endhead
1 & Welcome Back, Python & Values have types; f-strings build messages. Placement quiz (15 min). & data types, f-strings, \texttt{and}/\texttt{or}/\texttt{not} \\
2 & Functions That Give Back & \texttt{return} hands a value back; parameters make functions flexible. & \texttt{return}, docstrings, scope (light) \\
3 & Lists: Collections of Things & A list keeps many values in order. & lists, index, \texttt{append}, \texttt{len}, \texttt{for x in list} \\
4 & While Loops and Game Loops & Repeat until something changes; games are loops. & \texttt{while}, \texttt{break}, input validation, nested loops for pixel art \\
5 & Dictionaries and Secret Codes & Look things up by name, not number. & dict, keys/values, \texttt{.items()}, Caesar cipher \\
6 & Review and Checkpoint & Practice strategies for finding bugs. & tracebacks, print-debugging, rubber duck. \textbf{Checkpoint (15 min).} \\
7 & Files: Saving Your Stuff & Files let programs remember after they close. & \texttt{open}, \texttt{with}, read/write/append, \texttt{split}/\texttt{strip} \\
8 & Testing: Prove It Works & Tests are automatic checks that catch bugs early. & \texttt{assert}, pytest, normal and edge cases \\
9 & Event-Driven Games & Keys and timers make things happen; collisions end games. & \texttt{onkey}, \texttt{ontimer}, game state, \texttt{distance} collision. \textbf{Project kickoff.} \\
10 & Project Studio 1 & Plan, then build the core. & pseudocode, flowcharts, decomposition \\
11 & Project Studio 2 & Test, playtest, polish. & tests for own functions, feedback \\
12 & Showcase and Celebration & Share, explain, reflect. & explaining design. \textbf{Exit assessment (20 min).} \\ \midrule
B1 & Cipher \& Puzzle Day & Unplugged encryption and logic puzzles. & (use any week) \\
B2 & Generative Art Jam & Free-build with lists, random and functions. & (use any week) \\
\bottomrule
\end{longtable}}

%% ====================================================================
\section{4. Before you start: setup checklist}
Allow about one hour, once, before week 1. Check off each line.

\subsection{4.1 Thonny (on every Mac)}
\begin{checklist}
  \item Open Thonny and choose \menu{Help > About Thonny}. The Python version should be
        3.10 or newer. If it is older, download the current Mac version from
        \texttt{thonny.org} and drag it into Applications (replace the old one). No account is needed.
  \item If macOS says Thonny \enquote{cannot be opened}, right-click (or Control-click) Thonny
        in Applications and choose \menu{Open}, then \menu{Open} again.
  \item \menu{Tools > Options > Interpreter}: it should say \enquote{The same interpreter which runs
        Thonny (default)}. Leave it like that.
  \item Make the code easy to read: \menu{View > Increase font size} (or \keys{Cmd}\,+\,\keys{+}) until
        everyone at the screen can read it comfortably from where they sit.
\end{checklist}

\subsection{4.2 Install pytest (once per Mac, needed from week 8)}
pytest is a free helper that runs automatic tests. The \texttt{check\_my\_work.py} files in the
exercise folders also use it, so install it before week 1 to save time later.
\begin{enumerate}[itemsep=1pt]
  \item In Thonny choose \menu{Tools > Manage packages\ldots}
  \item Type \texttt{pytest} in the search box and click \menu{Search on PyPI}.
  \item Click \texttt{pytest} in the results, then \menu{Install}. Wait for it to finish (a minute or two).
  \item Close the window. Test it: open any \texttt{check\_my\_work.py} in an exercise folder and press
        \menu{Run}. You should see a short, friendly report rather than \enquote{I need a helper called pytest}.
\end{enumerate}
This needs the internet once. No account, login or personal data is involved.

\subsection{4.3 Copy the course folder to each Mac}
\begin{checklist}
  \item Copy the whole \texttt{kids-coding-curriculum} folder (USB stick, AirDrop or a shared drive) into each
        child's \menu{Documents} folder. Copy the \emph{whole} folder: the exercises use the
        \texttt{testing} folder inside it to check work.
  \item Level 2 lives in \texttt{kids-coding-curriculum/2-intermediate/}. The exercises are in
        \texttt{exercises/week01/}, \texttt{exercises/week02/} and so on (the spelling of
        \texttt{exercises} is intentional: it matches the original folder).
  \item Each exercise folder holds:
        \begin{itemize}[itemsep=0pt]
          \item \texttt{README.md}: the task, in kid-friendly words;
          \item \texttt{starter.py}: the file the child opens and edits;
          \item \texttt{check\_my\_work.py}: open it and press Run to check the work;
          \item \texttt{solution.py} and \texttt{test\_\ldots.py}: for you. If you prefer, delete
                \texttt{solution.py} from the children's copies.
        \end{itemize}
  \item Keep a clean copy of the folder on your own computer or a USB stick, so a broken file can
        always be replaced.
\end{checklist}

\subsection{4.4 Files and folders on a Mac (teach this in week 1, revisit in week 7)}
\begin{itemize}[itemsep=1pt]
  \item In \textbf{Finder}, show file extensions so \texttt{.py} and \texttt{.txt} are visible:
        \menu{Finder > Settings > Advanced}, check \enquote{Show all filename extensions}.
  \item In Thonny, open files with \menu{File > Open} and walk to
        \menu{Documents > kids-coding-curriculum > 2-intermediate > exercises > week\ldots}.
        Thonny remembers the last folder.
  \item Save new programs in the child's own folder with a clear name: lowercase, no spaces,
        ending in \texttt{.py} (for example \texttt{space\_snack.py}). Never name a file
        \texttt{turtle.py}, \texttt{random.py} or \texttt{test.py}: Python would load it instead of
        the real module and nothing works.
  \item A program reads and writes files \textbf{in the folder where the program is saved}. Always
        save a new program before pressing Run (see Troubleshooting).
\end{itemize}

\subsection{4.5 What to print}
\begin{tabularx}{\linewidth}{@{}>{\bfseries}p{4.8cm}X@{}}
\toprule
Print & When \\ \midrule
Assessments: placement quiz & Before week 1 (one per child). \\
Student Workbook chapter + homework set & Each week, one per child (or print the whole book at the start). \\
Cheat sheet (2 pages) & Week 1; one per child, keep beside the keyboard. Laminate if you can. \\
Progress tracker & Week 1; one per child plus your group page. \\
Parent notes & Week 1; one per family (including any visiting friend's family). \\
Assessments: checkpoint / exit & Week 6 / week 12. \\
Project brief & Week 9 (one per child). \\
Certificate & Week 12 (one per child). \\
Lesson plan, Answer Key pages & For you; reading on a screen is fine. \\
\bottomrule
\end{tabularx}

\subsection{4.6 Accessibility}
\begin{itemize}[itemsep=1pt]
  \item \textbf{Dyslexia-friendly edition:} every PDF also exists in a dyslexia-friendly version
        (Lexend font, larger text, wider spacing, cream paper color, no italics). After building
        the PDFs they are in \texttt{build/\allowbreak pdf/\allowbreak dyslexia/\allowbreak 2-intermediate/}. Offer it to any child
        who prefers it; it is not a label.
  \item \textbf{Thonny font size:} \menu{View > Increase font size}. For a lasting change use
        \menu{Tools > Options > Theme \& Font} and choose a larger editor and Shell font.
  \item \textbf{Color:} all course colors come from the Okabe--Ito palette, which works for the
        common kinds of color-blindness, and games always use shape and size as well as color.
        For extra help, macOS has color filters: \menu{System Settings > Accessibility > Display >
        Color Filters}. \menu{Zoom} on the same page magnifies the screen.
  \item \textbf{Keys:} if a child finds arrow keys awkward, any key can be used in a game
        (\texttt{"a"} and \texttt{"d"}, for example). Changing keys is a great extension task.
  \item \textbf{Reading aloud:} you may read any question, task or assessment aloud. We are
        assessing coding, not reading.
\end{itemize}

%% ====================================================================
\section{5. The 60-minute session}
Every lesson plan follows this template. Your prep is about 10 minutes: read the plan, open
the week's folder on each Mac, print the chapter and homework, set a timer.

\begin{timeplan}
  \timeslot{0--8}{Warm-up}{An unplugged activity or a quick predict-the-output puzzle. Collect last week's homework questions.}
  \timeslot{8--18}{Concept}{The one big idea of the day, with an analogy and the vocabulary box. Read the teacher script.}
  \timeslot{18--28}{Guided demo}{You (or a confident child) type the worked example while everyone predicts what will happen. Make a deliberate mistake and fix it together.}
  \timeslot{28--52}{Hands-on build}{Pairs work through the workbook exercises. Swap roles every 8 minutes. Bug hunt somewhere in this block. The last 5--6 minutes are creative free-build.}
  \timeslot{52--60}{Wrap-up}{Show-and-tell (one thing each), reflection prompt, exit ticket, hand out homework.}
\end{timeplan}

\begin{tip}[Timing rules of thumb]
\begin{itemize}[itemsep=1pt]
  \item Use a visible kitchen timer for the 8-minute swaps. The timer is the \enquote{bad guy}, not you.
  \item If the demo runs long, cut the second exercise, never the wrap-up: the exit ticket tells you
        what to do next week.
  \item Stop at 60 minutes even if the children want more. Leaving them wanting more is a feature.
\end{itemize}
\end{tip}

%% ====================================================================
\section{6. Pair programming at Level 2}
\begin{pairroles}[Three roles]
\textbf{Driver:} has the keyboard and mouse, types, and says out loud what they are typing.\\
\textbf{Navigator:} reads the task, thinks about the next step, spots typos, and says
\enquote{let's run it now}. The navigator does not touch the keyboard.\\
\textbf{Code reviewer (new in Level 2):} before each swap, the person who just navigated reads
the code for one minute and asks two questions from the review card:
\begin{itemize}[itemsep=0pt]
  \item \enquote{What will this line do?} (point at one line)
  \item \enquote{Is there a name that could be clearer?}
  \item \enquote{What happens if the number is 0, or the list is empty?} (from week 8)
  \item \enquote{What is one thing you did well?} (always end with this one)
\end{itemize}
\end{pairroles}

\begin{itemize}[itemsep=1pt]
  \item \textbf{4 children:} two pairs. Swap driver and navigator every 8 minutes; at each swap the
        navigator does a one-minute code review.
  \item \textbf{3 children:} a rotating trio: driver, navigator, reviewer. Rotate every 8 minutes so
        everyone does every role.
  \item \textbf{1 child:} the child drives; you navigate and review. Ask questions rather than giving
        instructions: \enquote{What should happen next?}
  \item \textbf{Very different levels in one pair:} the less confident child drives first. The
        navigator may only use words, never the mouse.
  \item \textbf{Fairness:} use the swap timer every time; praise good navigating as much as good typing.
\end{itemize}

%% ====================================================================
\section{7. The four weekly threads}
\begin{bughunt}[Every week]
A short broken program, usually based on the most common mistake of the day (for example
\texttt{"5" + 5}, an \texttt{IndexError}, a loop that never ends). Children predict what goes
wrong, run it, read the traceback, and fix it. From week 6 they name the strategy they used:
\textbf{read the traceback} (start at the last line), \textbf{print-debugging} (print the values),
or \textbf{rubber duck} (explain each line out loud to a toy or a partner).
Celebrate bugs: \enquote{Great, a bug! Now we get to be detectives.}
\end{bughunt}

\begin{unplugged}[Most weeks]
A paper or movement activity that makes an idea physical: type-sorting cards, function machines,
a human list, a paper Caesar wheel, trace tables, a game-loop board game. They need only paper,
pens and sticky notes. They are also your plan B if the computers misbehave.
\end{unplugged}

\begin{freebuild}[Last 5--6 minutes of the build]
The children change the day's program in any way they like: colors from the palette, new
shapes, a silly message, a different rule. Free-build is where many of them fall in love with
coding; protect these minutes.
\end{freebuild}

\begin{reflection}[Wrap-up]
One prompt each week, answered in the workbook or out loud: \enquote{What surprised you today?},
\enquote{Which bug took longest and how did you find it?}, \enquote{What would you teach a younger
child about lists?}. Short answers are fine. Read them: they are your best guide for next week.
\end{reflection}

%% ====================================================================
\section{8. Assessment plan}
Assessment in this course is \textbf{for learning}: it tells you what to do next. There are no
grades and no rankings. Results are shared kindly and privately.

\subsection{8.1 Placement quiz (start of week 1, 15 minutes)}
\begin{itemize}[itemsep=1pt]
  \item \textbf{What:} 8 short questions on Level 1 Python: predict the output, fix a bug, write a
        tiny loop. Paper and pencil, no computer. It is in the Assessments booklet.
  \item \textbf{When:} the first 15 minutes of week 1 (the week 1 plan allows for this), or a
        separate 15 minutes in the days before. Children new to the group (for example a friend
        joining) take it too.
  \item \textbf{Tell the children:} \enquote{This is not a test you pass or fail. It helps me choose
        the right challenges for you.}
  \item \textbf{Scoring:} use the Answer Key (Assessments section). It takes about 3 minutes per child.
  \item \textbf{Using the results:} the decision guide in the Answer Key places each child on the
        \textbf{support path}, the \textbf{standard path} or the \textbf{extension path}
        (see section~9). Write the path on the progress tracker. A child can change path at any
        time; look again after the week-6 checkpoint.
  \item \textbf{Pairs:} pair a support-path child with a standard-path child rather than with an
        extension-path child; the gap is easier to bridge.
\end{itemize}

\subsection{8.2 Formative checks (every week)}
Every lesson plan ends with a \textbf{formative check}: what to look for while the children work,
plus one exit-ticket question with its answer. Record a quick check mark, a dash or a question mark on
the progress tracker. Two question marks in a row on the same goal means: revisit it in the next
warm-up or use a buffer session.

\subsection{8.3 Mid-level checkpoint (week 6, about 15 minutes)}
Covers types, functions with \texttt{return}, lists, \texttt{while} loops and dictionaries:
predict, fix and write questions. Score it with the Answer Key, fill in the \enquote{I can} rubric
together with each child, and adjust paths for weeks 7--12.

\subsection{8.4 Exit assessment (week 12, about 20 minutes)}
Covers files, tests and event-driven programs, with read, predict, fix, write and explain
questions. Use it with the project rubric to fill in the final column of the progress tracker
and to write a sentence on each certificate.

\subsection{8.5 Final project rubric}
The project is assessed on four equal criteria, each on the four-level scale
(Just starting / Getting there / Got it / Expert):
\begin{enumerate}[itemsep=1pt]
  \item \textbf{Creativity:} the idea is their own and they made it fun or beautiful.
  \item \textbf{Working code:} it runs, it does what the pitch promised, it uses the must-haves.
  \item \textbf{Code quality and tests:} clear names, small functions, at least 3 passing tests.
  \item \textbf{Explaining it:} they can explain how one part works and one decision they made.
\end{enumerate}
The full rubric and scoring are in the Project Teacher Guide. A wonderful idea that is only half
working can still score highly: we value creativity and explanation as much as finished code.

\subsection{8.6 Progress tracker}
The progress tracker (Extras) has a page for each child (weekly log with \enquote{I can} statements)
and a group page for you (goals L2.G1--G9 at week 6 and week 12). Fill it in during the last
two minutes of the session, while the children write their reflection.

%% ====================================================================
\section{9. Differentiation: support and extension paths}
Every lesson plan has a support path and an extension path. Here is the whole level on one page so
you can plan ahead.

{\small
\begin{longtable}{@{}>{\bfseries}p{0.8cm}>{\raggedright\arraybackslash}p{6.9cm}>{\raggedright\arraybackslash}p{6.9cm}@{}}
\toprule
Wk & \faArrowDown\ Support path & \faArrowUp\ Extension path \\ \midrule\endhead
1 & Keep the type-sorting cards beside the keyboard; give an f-string with blanks to fill. & Format decimals with \texttt{f"\{x:.2f\}"}; a quiz that uses \texttt{and}/\texttt{or}. \\
2 & Use the paper function machine first; fill-in-the-blank \texttt{return} lines. & Functions that return True/False; a function that calls another function. \\
3 & Number the cards from 0 under the human list; use \texttt{for x in list} before indexes. & Build a list with a loop, then draw a palette stripe from it; \texttt{random.shuffle}. \\
4 & Trace table with the first rows filled in; start from a countdown. & Guessing game with a limited number of tries and \texttt{break}; pixel art from a list of strings. \\
5 & Paper Caesar wheel beside the keyboard; dictionaries with three keys only. & A decode function; count letters with a dictionary. \\
6 & Read the checkpoint aloud; focus on the predict questions. & Design a bug-hunt station for the others. \\
7 & Copy the \texttt{with open(\ldots)} lines from the cheat sheet; write first, read later. & Append to a log file; load a list of scores and turn them into numbers. \\
8 & Write tests for a given function, one \texttt{assert} at a time. & Edge-case tests; write a function to make given tests pass. \\
9 & Start from the game template with keys already bound; add one key. & A list of falling sprites; speed up over time. \\
10 & Starter template; aim for milestone M2 (core loop) only; you navigate. & Separate rules into a logic file; add levels or a menu. \\
11 & Adapt the example tests to their functions; polish one thing. & More tests (including a file test); let players choose their keys. \\
12 & Demo script card; show one feature; exit assessment read aloud. & Explain a design decision; help a younger sibling play. \\
B1 & Paper cipher wheel with a partner. & Invent a cipher with a dictionary. \\
B2 & Remix a past art exercise. & Generate art from a list of dictionaries; save a gallery. \\
\bottomrule
\end{longtable}}

%% ====================================================================
\section{10. The digital-citizenship thread}
Each week has a short \textbf{safety} conversation (2--4 minutes), usually in the wrap-up. They add
up to a real understanding of privacy, fairness and credit.

{\small
\begin{longtable}{@{}>{\bfseries}p{0.8cm}>{\raggedright\arraybackslash}p{4.4cm}>{\raggedright\arraybackslash}p{9.4cm}@{}}
\toprule
Wk & Theme & The key message \\ \midrule\endhead
1 & Hero names, not real names & Personal information is anything that identifies you: name, address, school, birthday, photos. Programs use a hero name. \\
2 & Comments for future-you & Code is read by people. Clear comments are a kindness to friends and to yourself next month. \\
3 & Color-blind-safe palettes & About 1 in 12 boys and 1 in 200 girls see some colors differently. Design for everyone: safe colors plus shapes and labels. \\
4 & Fair play & Good games do not trick or pressure players (no fake timers, no \enquote{pay to win}). \\
5 & Encryption and passwords & Secrets need protecting. Never share passwords; a Caesar cipher is fun but easy to break. \\
6 & Growth mindset & Mistakes are how programmers learn. Compare yourself with last month's you. \\
7 & What should (not) be saved & Only save what you need. Never save other people's private information in a file. \\
8 & Why apps need testing & Real bugs have cost money and caused problems; testing is how we earn people's trust. \\
9 & Attribution and open licenses & Use your own art or library art; credit other people's work; check the license. \\
10 & Splitting tasks fairly & In a team everyone gets interesting work and everyone's ideas are heard. \\
11 & Accessibility review & Check your own game: colors, shapes, text size, key choices. \\
12 & Showcase safety & Share with family only; no personal information on screen; nothing uploaded. \\
\bottomrule
\end{longtable}}

\begin{safety}[Our promises (repeat in week 1)]
No accounts, no logins, no uploading. Programs never ask for real names, ages, addresses,
schools, photos or passwords. Projects stay on our computers and are shown to family only.
\end{safety}

%% ====================================================================
\section{11. Troubleshooting}
{\small
\begin{longtable}{@{}>{\raggedright\arraybackslash}p{4.8cm}>{\raggedright\arraybackslash}p{9.8cm}@{}}
\toprule
\textbf{Problem} & \textbf{What to do} \\ \midrule\endhead
\multicolumn{2}{@{}l}{\textbf{Thonny}}\\ \midrule
Thonny will not open (\enquote{unidentified developer}) & Control-click Thonny in Applications, choose \menu{Open}, then \menu{Open}. Or \menu{System Settings > Privacy \& Security > Open Anyway}. \\
Nothing happens when pressing F5 & On many Macs F5 needs the \keys{fn} key. Use the green \menu{Run} button instead. \\
\texttt{IndentationError} or red underline & Thonny uses 4 spaces. Select the lines and use \menu{Edit > Indent} / \menu{Dedent}; never mix tabs and spaces. \\
The Shell shows a blinking cursor and nothing happens & The program is waiting for \texttt{input()}. Click in the Shell, type the answer, press \keys{return}. \\
\texttt{AttributeError: module 'turtle' has no attribute 'Turtle'} & A file in the same folder is called \texttt{turtle.py} (or \texttt{random.py}). Rename it, and delete the \texttt{\_\_pycache\_\_} folder next to it. \\ \midrule
\multicolumn{2}{@{}l}{\textbf{Turtle windows}}\\ \midrule
The turtle window does not appear & It is often behind Thonny. Look for the Python rocket icon in the Dock, or press \keys{Cmd}+\keys{Tab}. \\
The window closes at once or freezes with a spinning wheel & The program needs \texttt{turtle.done()} as its last line. Press the red \menu{Stop} button, add it, run again. \\
\texttt{turtle.Terminator} error after closing the window & Harmless: the program was still drawing when the window closed. Press \menu{Stop}, then \menu{Run}. \\
The game is slow or flickers & Use \texttt{screen.tracer(0)} at the start and \texttt{screen.update()} after drawing each frame (week 9). \\ \midrule
\multicolumn{2}{@{}l}{\textbf{Keys (\texttt{onkey}) do nothing}}\\ \midrule
No reaction at all & 1) Click the turtle window once so it has the keyboard focus. 2) Check that \texttt{screen.listen()} is called. 3) Check the key name: \texttt{"Left"}, \texttt{"Right"}, \texttt{"Up"}, \texttt{"Down"}, \texttt{"space"}, \texttt{"a"} (capital L for Left, small s for space). \\
The function runs once at the start, then never & \texttt{onkey(jump(), "space")} calls the function straight away. Remove the brackets: \texttt{onkey(jump, "space")}. \\ \midrule
\multicolumn{2}{@{}l}{\textbf{Files}}\\ \midrule
\texttt{FileNotFoundError} & The file does not exist in the folder where the program is saved. Check the spelling and extension (Finder may hide \texttt{.txt}). For high scores, catch it: \texttt{try} \ldots\ \texttt{except FileNotFoundError:} start with an empty list. \\
The file is created somewhere strange & The program was never saved (an \enquote{untitled} program runs from the home folder). Save it first. Programs in this course use \texttt{Path(\_\_file\_\_).parent} to keep files next to the program. \\
All my saved data disappeared & \texttt{open(name, "w")} empties the file first. Use \texttt{"a"} to add to the end. \\ \midrule
\multicolumn{2}{@{}l}{\textbf{pytest}}\\ \midrule
\enquote{No module named pytest} or \enquote{I need a helper called pytest} & Install it: \menu{Tools > Manage packages\ldots}, search \texttt{pytest}, \menu{Install} (section 4.2). Then restart Thonny. \\
Still not found after installing & \menu{Tools > Options > Interpreter} must be \enquote{The same interpreter which runs Thonny (default)}. Install again after changing it. \\
A test says \texttt{ModuleNotFoundError} for the child's program & The test file and the program must be in the same folder, and the \texttt{import} line must match the file name without \texttt{.py}. \\ \midrule
\multicolumn{2}{@{}l}{\textbf{Programs that never stop}}\\ \midrule
Infinite loop: nothing happens, or the Shell fills up & Press the red \menu{Stop} button (or \keys{Ctrl}+\keys{C} in the Shell). Then find the \texttt{while} loop and check that something inside it changes so the condition can become False. A great bug hunt! \\
\bottomrule
\end{longtable}}

%% ====================================================================
\section{12. Pacing, buffers and holidays}
Level 2 runs \textbf{10--12 weeks}, one session a week; weeks with no session (holidays, illness) do
not count. Use the \textbf{12-week plan} (the 12 core sessions in order) or the \textbf{10-week plan}
(merge weeks 4 + 5 and weeks 10 + 11, see 12.2).

\subsection{12.1 Golden rules}
\begin{itemize}[itemsep=1pt]
  \item \textbf{Never skip weeks 6 and 12.} Week 6 is the checkpoint that tells you whether the
        second half will work; week 12 is the showcase the children are working towards.
  \item The two \textbf{buffer sessions} (B1 Cipher \& Puzzle Day, B2 Generative Art Jam) are
        \textbf{optional}. Use them only when there is room inside the 12 weeks, for example in a
        10-week plan, before a holiday or when energy is low. They need almost no preparation.
  \item Missed a week? Do not double up and do not add weeks at the end: teach the missed week next
        time and merge two later weeks (12.2) so the level still ends within 12 weeks.
  \item Keep the project weeks (9--12) close together. A long break in the middle of the project
        is the hardest gap to recover from.
  \item After any break of two weeks or more, spend the warm-up on a 10-minute recap: a
        predict-the-output puzzle from the last lesson and one minute of \enquote{tell me what you
        remember}.
\end{itemize}

\subsection{12.2 Merging two weeks (10-week plan, or after lost sessions)}
In a merged session do both warm-ups as one, teach both concepts, and do only the first exercise of
each week; move the rest to homework.
\begin{itemize}[itemsep=1pt]
  \item \textbf{Weeks 4 + 5} (while and dictionaries) merge well; skip the pixel-art extension.
  \item \textbf{Weeks 10 + 11} merge for a smaller project (milestones M1--M3 only).
  \item \textbf{Weeks 1 + 2} (types and return) can merge instead of 4 + 5 for confident groups.
  \item Do \textbf{not} merge weeks 7 and 8 (files and testing are both new and both needed for
        the project), and never merge week 6 or week 12 with anything.
\end{itemize}

\subsection{12.3 Sample calendar: 12-week plan, summer 2027}
Twelve sessions on Saturdays, skipping the holiday weekends. If you add a summer break, put it
\emph{before} week 9, not in the middle of the project.

{\small
\begin{longtable}{@{}>{\bfseries}p{2.6cm}>{\raggedright\arraybackslash}p{4.8cm}>{\raggedright\arraybackslash}p{7.0cm}@{}}
\toprule
Date (2027) & Session & Notes \\ \midrule\endhead
Sat 8 May & Week 1: Welcome Back, Python & Placement quiz in the first 15 minutes. Hand out parent notes. \\
Sat 15 May & Week 2: Functions That Give Back & \\
Sat 22 May & Week 3: Lists & \\
Sat 29 May & --- (no session) & Memorial Day weekend. \\
Sat 5 Jun & Week 4: While Loops & \\
Sat 12 Jun & Week 5: Dictionaries & \\
Sat 19 Jun & --- (no session) & Juneteenth. \\
Sat 26 Jun & Week 6: Review and Checkpoint & Never skip. Update paths after the checkpoint. \\
Sat 3 Jul & --- (no session) & July 4 weekend. Optional: B1 puzzles as a car-journey activity. \\
Sat 10 Jul & Week 7: Files & 10-minute recap warm-up after the break. \\
Sat 17 Jul & Week 8: Testing & \\
Sat 24 Jul & Week 9: Event-Driven Games & Project kickoff. \\
Sat 31 Jul & Week 10: Project Studio 1 & Milestone M1 due at the start, M2 by the end. \\
Sat 7 Aug & Week 11: Project Studio 2 & M3 and M4. \\
Sat 14 Aug & Week 12: Showcase and Celebration & Never skip. Invite the families. Exit assessment. \\
\bottomrule
\end{longtable}}

\begin{teachernote}[Adapting the calendar]
For the 10-week plan, merge weeks 4 + 5 and weeks 10 + 11: the level then ends on 31 July, and you may
use B1 or B2 in the two free weeks if you like. Any weekday works just as well as Saturday.
\end{teachernote}

%% ====================================================================
\section{13. Coverage}
\subsection{13.1 Goals across the weeks}
\textbf{I} = introduced, R = revisited, \textbf{E}\textsuperscript{+} = extended. Marks in weeks 1--5 for
L2.G6 and L2.G8 revisit the debugging and pair work from Level 1.

{\small
\setlength{\tabcolsep}{3.2pt}
\begin{longtable}{@{}>{\bfseries}l>{\raggedright\arraybackslash}p{3.3cm}*{12}{c}cc@{}}
\toprule
Goal & Short name & 1 & 2 & 3 & 4 & 5 & 6 & 7 & 8 & 9 & 10 & 11 & 12 & B1 & B2 \\ \midrule\endhead
L2.G1 & types, f-strings & \cvI & \cvR & \cvR & & \cvR & \cvR & \cvE & & \cvR & \cvR & \cvR & & & \\
L2.G2 & while, nested, and/or & \cvI & & \cvR & \cvI & \cvR & \cvR & \cvR & & \cvE & \cvR & \cvR & & \cvR & \cvR \\
L2.G3 & functions, return & \cvR & \cvI & \cvR & \cvR & \cvR & \cvR & \cvR & \cvE & \cvR & \cvE & \cvR & \cvR & & \cvR \\
L2.G4 & lists, dictionaries & & & \cvI & \cvR & \cvE & \cvR & \cvR & \cvR & \cvE & \cvR & \cvR & & \cvR & \cvR \\
L2.G5 & files & & & & & & & \cvI & \cvR & & \cvR & \cvE & & & \cvR \\
L2.G6 & debugging, tests & \cvR & \cvR & \cvR & \cvR & \cvR & \cvI & \cvR & \cvI & \cvR & \cvR & \cvE & \cvR & \cvR & \\
L2.G7 & events, games, art & & & \cvR & \cvE & & & & & \cvI & \cvE & \cvE & \cvR & & \cvE \\
L2.G8 & plan, team, feedback & \cvR & \cvR & \cvR & \cvR & \cvR & \cvR & \cvR & \cvR & \cvR & \cvI & \cvE & \cvR & \cvR & \cvR \\
L2.G9 & safe, fair, private & \cvI & \cvR & \cvR & \cvR & \cvE & & \cvE & \cvR & \cvR & \cvR & \cvR & \cvR & \cvR & \\
\bottomrule
\end{longtable}}

\subsection{13.2 Objectives and the lessons that teach them}
The table below is generated automatically from the lesson plans: it lists every learning
objective and the lesson where it is taught.

\kcinput{generated/coverage}

%% ====================================================================
\section{14. Spiral map: from Level 1 to Level 3}
Big ideas come back again and again, each time a little deeper. Use this map to connect today's
lesson with what the children already know (\enquote{Remember My Blocks in Scratch? A function
with \texttt{return} is the same idea, grown up.}) and to see where it leads.

{\small
\begin{longtable}{@{}>{\bfseries\raggedright\arraybackslash}p{2.8cm}>{\raggedright\arraybackslash}p{3.7cm}>{\raggedright\arraybackslash}p{4.1cm}>{\raggedright\arraybackslash}p{3.9cm}@{}}
\toprule
Big idea & Level 1 (what they know) & Level 2 (this level) & Level 3 (where it goes) \\ \midrule\endhead
Sequence \& algorithms & W1 introduced (unplugged), W7 revisited in text & W2 functions, W10 pseudocode & W4--5 named algorithms, efficiency \\
Loops & W2 repeat in Scratch, W8 \texttt{for} & W4 \texttt{while} and nested loops, W9 game loop & W4 search loops, W7 recursion \\
Events \& conditionals & W3 events, W9 \texttt{if}/\texttt{elif} & W1 \texttt{and}/\texttt{or}, W9 \texttt{onkey}/\texttt{ontimer} & W2--3 methods that react, project \\
Variables \& data & W4 variables, W9 \texttt{input} & W1 types, W3 lists, W5 dictionaries, W7 files & W2 attributes, W5 stacks/queues, W7 grids, W8 CSV data \\
Abstraction / functions & W5 My Blocks, W10 \texttt{def} & W2 \texttt{return}, W8 testable functions & W2--3 classes, W10 modules \\
Debugging \& testing & weekly bug hunts, W6 relay & W6 strategies, W8 pytest & W6 debugger, project with 5+ tests \\
Collaboration & driver/navigator every week & code reviewer role, W10 task split & W1 git, W10 code review \\
Art \& games & W2 pen art, W4 catch game, W10 random art & W3 palettes, W4 pixel art, W9 games & W3 shape families, W7 fractals, project \\
Citizenship \& ethics & privacy, credit, kindness, accessibility & data privacy, encryption, fair games, attribution & git secrets, consent, bias, AI ethics, licenses \\
\bottomrule
\end{longtable}}

\begin{tip}[Looking ahead to Level 3]
Level 3 starts with a placement quiz on Level 2 ideas (functions, lists, dictionaries, files,
tests). Children who finish Level 2 comfortable with the exit assessment are ready. Keep their
final projects: Level 3 opens by putting old code under version control with git.
\end{tip}

\end{document}
